\subsection{衰变链的相等常数极限}
一般两级模型中令 $\lambda_2\to\lambda_1=\lambda$，巴特曼式表面上为 $0/0$，应先取极限而非宣布发散。用导数或直接积分得
\[
 N_2(t)=\lambda N_{10}t e^{-\lambda t},\qquad
 t_{\max}=\frac1\lambda,\qquad N_{2,\max}=\frac{N_{10}}e.
\]
无论两常数是否相等，子核峰值时都满足生产率等于损失率：$\lambda_1N_1=\lambda_2N_2$。相等常数极限、初始导数 $\dot N_2(0)=\lambda_1N_{10}$、长时间趋零，是检查衰变链公式的三个独立办法。若母核只有比例 $b$ 的分支通向所计子核，源项应改为 $b\lambda_1N_1$；若初始已有子核，还需加 $N_{20}e^{-\lambda_2t}$。
